% Volume IX: Integration, Capstones, and Reference Material
% Continuity note for Chapter 53.
% Previous dependencies: Chapters 5--7 provide least squares, PCA, matrix derivatives, and backpropagation shapes; Chapters 13--18 provide probability, Bayes, and information; Chapters 22, 27, 31--42 provide variational objectives, stochastic processes, learning, deep networks, and generative models; Chapters 47--52 provide neural coding, manifolds, predictive coding, neural data analysis, and Bio-AI comparison.
% New durable concepts: integrated mathematical case-study reasoning; linear regression as the base case for deep networks; PCA as the base case for neural manifold analysis; Bayes as the base case for predictive-coding updates; stochastic processes as the base case for diffusion models and neural noise.
% Forward hooks: Chapter 54 collects the formulas, notation, derivations, and computational recipes used in these case studies.
% Notation commitments: Data matrices are X; responses or targets are Y; regression weights are W; neural population activity is R or X; latent coordinates are Z; hidden causes are z; observations are y; Brownian motion is W_t; stochastic states are X_t.
\section{Chapter 53: Integrated Mathematical Case Studies}
\label{ch:integrated-mathematical-case-studies}

The course has introduced many mathematical tools: linear algebra, calculus, probability, information, optimization, stochastic processes, machine learning, neural data analysis, and circuit models. The purpose of this chapter is to show how those tools combine in complete modeling stories. Each section begins with a familiar mathematical object and extends it toward a modern neuroscience or AI application.

\paragraph{Chapter problem.}
How do the course's core mathematical ideas reappear, with additional structure, in deep networks, neural manifolds, predictive coding, diffusion models, and neural noise?

\paragraph{Section map.}
Section 53.1 connects linear regression to deep networks through composition, loss minimization, and backpropagation. Section 53.2 connects PCA to neural manifolds through covariance, projections, latent coordinates, and information loss. Section 53.3 connects Bayes' theorem to predictive coding through Gaussian latent-variable models and precision-weighted prediction errors. Section 53.4 connects stochastic processes to diffusion models and neural noise through Brownian increments, SDEs, OU dynamics, and denoising.

\subsection{53.1 From linear regression to deep networks}

\paragraph{Guiding question.}
In what precise sense is a deep network a nonlinear, compositional extension of linear regression?

\paragraph{Beginner-facing intuition.}
Linear regression learns weights that map input features to outputs. A deep network also learns weights, but it inserts nonlinear hidden representations between input and output. The basic ingredients remain recognizable: a parameterized function, a loss, gradients, train/test evaluation, and shape discipline. Deep learning becomes less mysterious when it is seen as regression with learned intermediate features and repeated chain rules.

\paragraph{Formal setup.}
Let \(X\in\mathbb R^{M\times d}\) contain \(M\) examples and \(d\) input features. Let \(Y\in\mathbb R^{M\times k}\) contain targets. Multivariate linear regression fits
\[
\widehat Y=XW,\qquad W\in\mathbb R^{d\times k},
\]
by minimizing
\[
L(W)=\frac{1}{2}\|XW-Y\|_F^2.
\]
A one-hidden-layer neural network replaces the direct linear map by
\[
H=\sigma(XW_1+\mathbf 1 b_1^\top),\qquad
\widehat Y=HW_2+\mathbf 1 b_2^\top,
\]
where \(W_1\in\mathbb R^{d\times h}\), \(W_2\in\mathbb R^{h\times k}\), and \(\sigma\) is applied elementwise. A deep network repeats this pattern:
\[
h_0=x,\qquad h_{\ell+1}=\sigma_\ell(W_\ell h_\ell+b_\ell),\qquad f_\theta(x)=h_L.
\]

\paragraph{Core derivation: regression gradient to backpropagation.}
For linear regression,
\[
dL=\langle XW-Y, X\,dW\rangle_F
=\langle X^\top(XW-Y),dW\rangle_F.
\]
Thus
\[
\boxed{\nabla_W L=X^\top(XW-Y).}
\]
Setting the gradient to zero gives the normal equations \(X^\top XW=X^\top Y\).

For the one-hidden-layer network with squared loss, define
\[
E=\widehat Y-Y.
\]
The output-layer gradient is
\[
\nabla_{W_2}L=H^\top E.
\]
The hidden error is
\[
\Delta_H=(EW_2^\top)\odot \sigma'(XW_1+\mathbf 1 b_1^\top),
\]
so
\[
\nabla_{W_1}L=X^\top \Delta_H.
\]
The pattern is the same as linear regression: residuals are pushed backward through transposed linear maps and multiplied by local derivatives. Deep backpropagation repeats this mechanism layer by layer.

\paragraph{Concrete example.}
Let \(x\in\mathbb R\), \(y\in\mathbb R\), and consider linear regression \(\widehat y=wx\). For one example \(x=2,y=5,w=1\),
\[
L=\frac12(wx-y)^2=\frac12(2-5)^2=4.5,
\]
and
\[
\frac{dL}{dw}=(wx-y)x=(-3)(2)=-6.
\]
A gradient step with learning rate \(0.1\) gives \(w_{\mathrm{new}}=1-0.1(-6)=1.6\). A deep network uses the same loss-gradient-update logic, but the derivative is routed through hidden layers.

\paragraph{Computational neuroscience example.}
A linear encoding model predicts a neuron's response from stimulus features. A deep encoding model replaces hand-chosen features with activations from several learned or pretrained layers. The scientific question changes from ``which linear feature predicts the neuron?'' to ``which learned representation, mapped through a constrained readout, predicts the response on held-out stimuli?''

\paragraph{AI or machine-learning example.}
In image classification, the last layer of a neural network is often a linear classifier on learned features. Training the whole network jointly learns both the features and the linear readout. Fine-tuning only the last layer is close to ordinary regression or classification on fixed features.

\paragraph{Common misconceptions and failure modes.}
Deep networks are not just bigger linear models: without nonlinearities, a product of linear layers collapses to one linear map. More layers increase expressivity but also introduce nonconvex optimization and identifiability issues. A low training loss does not imply causal explanation or biological plausibility. Gradient formulas require shape discipline; transposes and elementwise derivatives are not optional decorations.

\paragraph{Exercises.}
\begin{enumerate}[label=\textbf{53.1.\arabic*.}, leftmargin=*]
\item \textbf{Mechanical.} State the shapes of \(X\), \(W\), \(Y\), and \(XW-Y\) when \(M=20\), \(d=4\), and \(k=3\).
\item \textbf{Proof or derivation.} Derive \(\nabla_W \frac12\|XW-Y\|_F^2=X^\top(XW-Y)\) using differentials.
\item \textbf{Numerical/computational.} For \(x=3\), \(y=7\), and \(w=2\), compute squared-error loss \(\frac12(wx-y)^2\) and gradient \((wx-y)x\).
\item \textbf{Interpretation.} Explain why a deep network with no nonlinearities between layers is equivalent to a single linear map.
\item \textbf{Misconception/debugging.} A student says, ``Backpropagation is unrelated to regression.'' Correct the statement by comparing residual propagation in linear regression and a one-hidden-layer network.
\end{enumerate}

\subsection{53.2 From PCA to neural manifolds}

\paragraph{Guiding question.}
How does PCA become the starting point for modern analyses of low-dimensional neural manifolds?

\paragraph{Beginner-facing intuition.}
PCA finds directions of large variance in a cloud of points. In neural data, each point can be a population activity vector. If the cloud lies near a low-dimensional set, then a few coordinates may describe the dominant neural activity patterns. Neural manifold analysis extends this idea: the low-dimensional set may be curved, task-dependent, dynamic, and linked to behavior.

\paragraph{Formal setup.}
Let
\[
X\in\mathbb R^{M\times N}
\]
be a centered neural activity matrix with \(M\) observations and \(N\) neurons. The empirical covariance is
\[
C=\frac{1}{M-1}X^\top X.
\]
Let \(v_1,\ldots,v_N\) be orthonormal eigenvectors of \(C\) with eigenvalues
\[
\lambda_1\geq\cdots\geq\lambda_N\geq0.
\]
The \(q\)-dimensional PCA coordinates are
\[
Z=XV_q,\qquad V_q=[v_1\ \cdots\ v_q].
\]
A neural manifold model assumes that activity is concentrated near a low-dimensional set
\[
\mathcal M\subseteq\mathbb R^N
\]
with intrinsic dimension \(q\ll N\). Linear PCA uses a \(q\)-dimensional affine subspace as \(\mathcal M\). Nonlinear manifold methods allow curved \(\mathcal M\).

\paragraph{Core derivation: PCA as optimal linear reconstruction.}
For an orthonormal matrix \(V_q\in\mathbb R^{N\times q}\), the rank-\(q\) reconstruction is
\[
\widehat X=XV_qV_q^\top.
\]
The reconstruction error is
\[
\|X-\widehat X\|_F^2
=\|X\|_F^2-\|XV_q\|_F^2.
\]
Thus minimizing reconstruction error is equivalent to maximizing
\[
\|XV_q\|_F^2
=\operatorname{tr}(V_q^\top X^\top X V_q)
=(M-1)\operatorname{tr}(V_q^\top C V_q).
\]
The maximizing orthonormal columns are the top eigenvectors of \(C\). The explained variance fraction is
\[
\frac{\lambda_1+\cdots+\lambda_q}{\lambda_1+\cdots+\lambda_N}.
\]

\paragraph{Concrete example.}
Suppose the covariance eigenvalues of a neural population are
\[
9,\quad 4,\quad 1,\quad 1.
\]
The first two PCs explain
\[
\frac{9+4}{9+4+1+1}=\frac{13}{15}\approx86.7\%
\]
of the variance. A two-dimensional plot may therefore capture much of the activity variance, but it still discards about \(13.3\%\).

\paragraph{Information-theory connection.}
If \(S\) is a stimulus or behavior variable, \(X\) is full neural activity, and \(Z=XV_q\) is a deterministic projection, then
\[
S\to X\to Z
\]
is a Markov chain. By data processing,
\[
\boxed{I(S;Z)\leq I(S;X).}
\]
Projection cannot create information about \(S\) that was not present in the full activity. A low-dimensional representation can still be useful if it preserves most task-relevant information while discarding noise. This distinction matters: variance explained and information about behavior are related only under additional assumptions.

\paragraph{Computational neuroscience example.}
Motor-cortex activity during reaching often occupies a low-dimensional trajectory when projected with PCA or factor analysis. A manifold coordinate may correlate with movement direction, speed, or preparation. The claim should be checked with held-out behavior and not inferred solely from a pretty two-dimensional plot.

\paragraph{AI or machine-learning example.}
Hidden activations in a deep network also form representation clouds. PCA can reveal whether class labels, time, or task variables organize the activity. Nonlinear embeddings can visualize structure, but quantitative claims require held-out decoding, reconstruction, or information estimates.

\paragraph{Common misconceptions and failure modes.}
High variance is not automatically task relevance. A PC can capture movement artifact or arousal instead of computation. A 2D plot is a projection, not the whole manifold. Nonlinear embeddings such as t-SNE or UMAP can distort distances and densities. Information estimates from projected data can be biased by finite samples, binning, and decoder choice.

\paragraph{Exercises.}
\begin{enumerate}[label=\textbf{53.2.\arabic*.}, leftmargin=*]
\item \textbf{Mechanical.} If \(X\in\mathbb R^{200\times 50}\) and \(q=3\), what are the shapes of \(V_q\), \(Z\), and \(XV_qV_q^\top\)?
\item \textbf{Proof or derivation.} Derive why maximizing \(\operatorname{tr}(V_q^\top C V_q)\) over orthonormal \(V_q\) selects the top covariance eigenvectors.
\item \textbf{Numerical/computational.} Eigenvalues are \(5,3,2,0\). Compute the variance explained by the first one, first two, and first three PCs.
\item \textbf{Interpretation.} Explain why a low-dimensional projection can preserve behaviorally relevant information even though data processing says it cannot increase information.
\item \textbf{Misconception/debugging.} A student says, ``The first two PCs are the neural manifold.'' Correct the statement using linear projections and intrinsic geometry.
\end{enumerate}

\subsection{53.3 From Bayes to predictive coding}

\paragraph{Guiding question.}
How does Bayesian inference in a simple generative model turn into prediction-error minimization?

\paragraph{Beginner-facing intuition.}
Bayes starts with a latent cause and an observation. The generative model runs forward from cause to observation; inference runs backward from observation to cause. Predictive coding can be understood as an algorithmic proposal for approximate inference: maintain a current estimate of the latent cause, predict the observation, compute prediction errors, and update the estimate to reduce precision-weighted error.

This does not prove that every cortical circuit performs exact Bayesian inference. It shows how a particular mathematical objective gives prediction-error update equations that can be compared with neural signals.

\paragraph{Formal setup.}
Let a latent cause \(z\in\mathbb R^d\) have Gaussian prior
\[
z\sim\mathcal N(\mu_0,\Sigma_0),
\]
and let an observation \(y\in\mathbb R^p\) follow a linear-Gaussian likelihood
\[
y\mid z\sim\mathcal N(Cz,\Sigma_y).
\]
The posterior satisfies
\[
p(z\mid y)\propto p(y\mid z)p(z).
\]
Ignoring constants independent of \(z\), the negative log posterior is
\[
\mathcal F(z)
=\frac12(y-Cz)^\top\Sigma_y^{-1}(y-Cz)
+\frac12(z-\mu_0)^\top\Sigma_0^{-1}(z-\mu_0).
\]
This is a precision-weighted prediction-error objective. Define
\[
\varepsilon_y=y-Cz,\qquad \varepsilon_z=z-\mu_0,
\]
with precisions \(\Pi_y=\Sigma_y^{-1}\) and \(\Pi_0=\Sigma_0^{-1}\).

\paragraph{Core derivation: posterior mode and update.}
Differentiate \(\mathcal F(z)\):
\[
\nabla_z\mathcal F(z)
=-C^\top\Pi_y(y-Cz)+\Pi_0(z-\mu_0).
\]
Gradient descent on \(\mathcal F\) gives
\[
\dot z
=-\nabla_z\mathcal F(z)
=C^\top\Pi_y\varepsilon_y-\Pi_0\varepsilon_z.
\]
The update increases \(z\) in directions that reduce sensory prediction error and decreases it in directions that violate the prior. Setting the gradient to zero gives the Gaussian posterior mean and mode:
\[
(C^\top\Pi_y C+\Pi_0)z
=C^\top\Pi_y y+\Pi_0\mu_0.
\]

\paragraph{Concrete example.}
Let \(d=p=1\), \(C=1\), prior \(z\sim\mathcal N(0,1)\), and observation noise variance \(1\). For observed \(y=2\),
\[
\mathcal F(z)=\frac12(2-z)^2+\frac12 z^2.
\]
The derivative is
\[
\mathcal F'(z)=-(2-z)+z=2z-2.
\]
The optimum is \(z=1\), halfway between the prior mean \(0\) and observation \(2\) because the prior and likelihood have equal variance.

\paragraph{Computational neuroscience example.}
In a hierarchical sensory model, higher-level activity can be interpreted as a latent cause estimate, feedback as prediction, and feedforward signals as prediction errors. This mapping is a hypothesis about circuit organization. It becomes testable only when the variables, precisions, timing, and perturbation predictions are specified.

\paragraph{AI or machine-learning example.}
Autoencoders and variational autoencoders also use latent variables to reconstruct observations. A VAE optimizes a variational objective involving reconstruction error and KL divergence. Predictive-coding networks can be viewed as iterative inference networks that minimize local prediction errors instead of performing one feedforward pass.

\paragraph{Common misconceptions and failure modes.}
Prediction error is not automatically surprise in the information-theoretic sense unless a probability model defines it. Predictive coding is not identical to next-token prediction in language models, though both involve prediction losses. Precision weighting matters: a large error from an unreliable channel should not dominate a small error from a reliable channel. Biological interpretations require caution because many different circuits can implement similar update equations.

\paragraph{Exercises.}
\begin{enumerate}[label=\textbf{53.3.\arabic*.}, leftmargin=*]
\item \textbf{Mechanical.} Identify the sensory prediction error and prior prediction error in \(\mathcal F(z)\).
\item \textbf{Proof or derivation.} Differentiate \(\mathcal F(z)\) and derive the linear system for the posterior mode.
\item \textbf{Numerical/computational.} In the scalar model with prior variance \(4\), observation noise variance \(1\), prior mean \(0\), and observation \(y=10\), compute the posterior mode.
\item \textbf{Interpretation.} Explain why increasing sensory noise reduces the influence of \(y\) on the posterior estimate.
\item \textbf{Misconception/debugging.} A student says, ``Predictive coding proves the brain minimizes free energy.'' Rewrite this as a precise model-dependent claim.
\end{enumerate}

\subsection{53.4 From stochastic processes to diffusion models and neural noise}

\paragraph{Guiding question.}
How do stochastic-process tools describe both noisy neural dynamics and diffusion-based generative models?

\paragraph{Beginner-facing intuition.}
A stochastic process is a time-indexed collection of random variables. Neural activity fluctuates over time because of synaptic noise, spiking variability, unobserved inputs, and changing internal state. Diffusion generative models deliberately add noise to data over time and learn how to reverse that corruption. Both use the mathematics of random paths, transition kernels, Gaussian increments, and stochastic differential equations.

\paragraph{Formal setup.}
A stochastic process is
\[
\{X_t:t\in T\}
\]
on a probability space, where each \(X_t\) is a random variable. A continuous-time SDE has the form
\[
dX_t=f(X_t,t)\,dt+g(X_t,t)\,dW_t,
\]
where \(W_t\) is Brownian motion. Brownian increments satisfy
\[
W_{t+\Delta t}-W_t\sim\mathcal N(0,\Delta t)
\]
and are independent over disjoint intervals.

The Ornstein-Uhlenbeck process is
\[
dX_t=-\alpha(X_t-\mu)\,dt+\sigma\,dW_t,
\]
with mean reversion rate \(\alpha>0\) toward \(\mu\) and noise scale \(\sigma\geq0\). Its stationary distribution is
\[
X_\infty\sim \mathcal N\!\left(\mu,\frac{\sigma^2}{2\alpha}\right).
\]

A discrete diffusion model often uses a forward noising chain
\[
q(x_t\mid x_{t-1})=\mathcal N(\sqrt{1-\beta_t}\,x_{t-1},\beta_t I),
\]
where \(0<\beta_t<1\). The learned reverse model approximates \(p_\theta(x_{t-1}\mid x_t)\).

\paragraph{Core derivation: closed-form forward noising.}
Define \(\alpha_t=1-\beta_t\) and \(\bar\alpha_t=\prod_{s=1}^t\alpha_s\). The forward recursion is
\[
x_t=\sqrt{\alpha_t}x_{t-1}+\sqrt{1-\alpha_t}\,\epsilon_t,
\qquad \epsilon_t\sim\mathcal N(0,I).
\]
By repeatedly substituting and using independence of Gaussian noise, one obtains
\[
\boxed{
q(x_t\mid x_0)=
\mathcal N(\sqrt{\bar\alpha_t}\,x_0,(1-\bar\alpha_t)I).
}
\]
Equivalently,
\[
x_t=\sqrt{\bar\alpha_t}\,x_0+\sqrt{1-\bar\alpha_t}\,\epsilon,\qquad \epsilon\sim\mathcal N(0,I).
\]
This formula lets a diffusion model train at arbitrary noise levels without simulating every previous step.

\paragraph{Concrete example.}
Let \(x_0=10\), \(\beta_1=\beta_2=0.1\), so \(\alpha_1=\alpha_2=0.9\) and \(\bar\alpha_2=0.81\). Then
\[
x_2\mid x_0\sim\mathcal N(\sqrt{0.81}\cdot10,(1-0.81))
=\mathcal N(9,0.19)
\]
in one dimension. The mean has moved toward zero and the variance has increased.

\paragraph{Computational neuroscience example.}
The OU process is a simple model for fluctuating membrane potential around a mean input or for latent neural state with restoring dynamics. The drift \(-\alpha(X_t-\mu)\) prevents variance from growing without bound, while \(\sigma dW_t\) injects ongoing noise. This is more realistic than pure Brownian motion for many stable neural variables.

\paragraph{AI or machine-learning example.}
Score-based and diffusion generative models learn how data distributions change under noise and how to denoise. In continuous time, some formulations use forward SDEs and learned reverse-time SDEs or probability-flow ODEs. The same \(dW_t\) notation appears, but the scientific interpretation differs: in neural models, noise is part of the modeled biological variability; in diffusion models, noise is an algorithmic construction for generation.

\paragraph{Common misconceptions and failure modes.}
The symbol \(dW_t\) is not an ordinary small number; Brownian increments scale like \(\sqrt{dt}\), not \(dt\). Pure Brownian motion has variance that grows forever, while OU noise has a stationary variance. A diffusion generative model is not claiming that images are produced by physical Brownian motion. Neural variability is not automatically nuisance noise; it may reflect latent variables, computation, or unmeasured inputs.

\paragraph{Exercises.}
\begin{enumerate}[label=\textbf{53.4.\arabic*.}, leftmargin=*]
\item \textbf{Mechanical.} State the drift and diffusion terms in \(dX_t=-\alpha(X_t-\mu)\,dt+\sigma\,dW_t\).
\item \textbf{Proof or derivation.} Derive \(q(x_t\mid x_0)=\mathcal N(\sqrt{\bar\alpha_t}x_0,(1-\bar\alpha_t)I)\) by induction for the discrete forward diffusion chain.
\item \textbf{Numerical/computational.} If \(\bar\alpha_t=0.25\), \(x_0=4\), and \(\epsilon=-1\), compute \(x_t=\sqrt{\bar\alpha_t}x_0+\sqrt{1-\bar\alpha_t}\epsilon\) in one dimension.
\item \textbf{Interpretation.} Why is an OU process often a better model of stable neural fluctuations than Brownian motion?
\item \textbf{Misconception/debugging.} A student says, ``Diffusion models prove that neural noise is just image-generation noise.'' Correct the statement by distinguishing shared mathematics from domain interpretation.
\end{enumerate}

\paragraph{Closing bridge.}
These four case studies show how the course's mathematical foundations compose. Regression becomes deep learning when features are learned by differentiable composition. PCA becomes manifold analysis when population activity is treated geometrically and linked to information and behavior. Bayes becomes predictive coding when posterior inference is implemented by precision-weighted prediction-error descent. Stochastic processes become a shared language for neural variability and diffusion-based generation. Chapter 54 turns the formulas used here into compact reference tables and computational recipes.

\clearpage
